\section{随机过程、相关结构与气候形变理论}
\label{sec:scenario-stochastic-theory}

\begin{theorynote}
本章推导当前二元相关过程及 profile 乘法形变。核心区分是：确定性季节/日内基线描述条件均值，
随机倍率描述围绕该均值的离散不确定性；气候形变改变条件均值，不能与随机噪声混为一项。
\end{theorynote}

\subsection{乘法分解}
对资产族 $a\in\{L,PV,W,R\}$，场景功率写为
\begin{equation}
 P^a_t=P^{a,0}\,s^a_t\,c^a_t(\mathrm{SSP},y)\,m^a_t\,h^a_t,
\end{equation}
其中 $P^{a,0}$ 为资产基准，$s_t$ 为确定性日内/季节形状，$c_t$ 为气候形变，$m_t$ 为随机
倍率，$h_t$ 为台风等事件影响。未启用某层时对应因子取 1。PV 的夜间 $s^{PV}_t=0$，因此
不能通过 $P_t/P^0_t$ 在夜间恢复随机倍率。

\subsection{平稳二元 AR(1)}
令 $x_t=(x^L_t,x^R_t)^\mathsf T$，同一时间相关参数作用于两维：
\begin{equation}
 x_t=\rho_t x_{t-1}+\sqrt{1-\rho_t^2}\,\varepsilon_t,
 \qquad |\rho_t|<1.
\end{equation}
创新由两个独立标准正态变量 $u_t,v_t$ 构造：
\begin{equation}
 \varepsilon^L_t=u_t,\qquad
 \varepsilon^R_t=\rho_{LR}u_t+\sqrt{1-\rho_{LR}^2}\,v_t.
\end{equation}
所以 $\operatorname{Var}(\varepsilon^L)=\operatorname{Var}(\varepsilon^R)=1$，且
$\operatorname{Corr}(\varepsilon^L,\varepsilon^R)=\rho_{LR}$。若初始状态也服从同一二元标准
正态，则
\begin{align}
 \operatorname{Var}(x_t)&=\rho_t^2 I+(1-\rho_t^2)I=I,\\
 \operatorname{Cov}(x_t,x_{t-\ell})&=\rho_t^\ell I,\\
 \operatorname{Corr}(x^L_t,x^R_t)&=\rho_{LR}.
\end{align}
这给出实现所需的目标时间相关和同期交叉相关。

\subsection{块结点与线性插值}
若 \fld{block\_hours} 为 $b>1$，AR(1) 在结点 $j$ 上生成，再对 $t=jb+r$ 做
\begin{equation}
 \widetilde x_t=(1-r/b)x_j+(r/b)x_{j+1}.
\end{equation}
插值改变逐时边际方差和 lag-1 相关；因此 $\rho_t$ 是结点过程参数，不保证插值后每小时样本
相关恰等于输入值。数值相关性验证使用 $b=1$，以隔离 AR(1) 本身。

\subsection{倍率截断与矩偏差}
实现采用仿射映射和硬截断
\begin{align}
 m^L_t&=\operatorname{clip}(1+\sigma_Lx^L_t,l_L,u_L),\\
 m^R_t&=\operatorname{clip}(1+\sigma_Rx^R_t,l_R,u_R).
\end{align}
未触发截断时，$E[m]=1$、$\operatorname{sd}(m)=\sigma$；截断后均值、方差和相关性均发生
偏移。截断概率可由标准正态 CDF 估计：
\begin{equation}
 P_{clip}=\Phi\!\left(\frac{l-1}{\sigma}\right)
 +1-\Phi\!\left(\frac{u-1}{\sigma}\right).
\end{equation}
因此统计测试必须同时报告 multiplier 的最小、最大、均值和样本标准差，不能只检查相关系数。

\subsection{PV 周期性选择偏差}
若仅在 $s^{PV}_t>0$ 时用 $P^{PV}_t/P^{PV,base}_t$ 恢复新能源倍率，观测索引集合
$D=\{t:s^{PV}_t>0\}$ 是确定性的周期性子样本。对自相关过程，
\begin{equation}
 \widehat\rho_D=\operatorname{Corr}(x^L_t,x^R_t\mid t\in D)
\end{equation}
仍以 $\rho_{LR}$ 为总体目标，但有效样本量和有限样本偏差不同。固定算例实测全天风电相关
-0.290252，而 PV 白天子样本为 -0.341279；后者只作诊断，不作为 AR(1) 准入门。

\subsection{储能 SOC 扰动}
对储能设备 $j$，先在 multiplier 空间做最多 128 次截断正态拒绝抽样，再回退 clip：
\begin{equation}
 \widetilde s_j=\operatorname{clip}(s^0_jm_j,s_j^{min},s_j^{max}).
\end{equation}
聚合能量和 SOC 为
\begin{equation}
 E^{st}=\sum_j E_j^{rated}\widetilde s_j,\qquad
 s^{agg}=\frac{E^{st}}{\sum_jE_j^{rated}}.
\end{equation}
当前生成器把该值填充到整个候选时域，不模拟充放电动态；动态 SOC 必须由 time-series
调度/PF 流水线计算。

\subsection{气候形变与随机气候扰动}
月度温升 $\Delta T_m$ 和辐照变化 $\Delta GHI_m$ 来自资源查找；缺失时采用零增量并记录
\fld{climate\_data\_found=false}。对负荷，可概括为
\begin{equation}
 c^L_t=1+\beta_T\Delta T_{m(t)}+\eta^L_{m(t)},
\end{equation}
对 PV 则先改变辐照参考，再施加温度效率项。随机气候扰动使用独立
\fld{climate\_seed}；它与二元 AR(1) 的 seed 空间分离。

\begin{gapnote}
当前模型没有空间高斯场、copula、多变量 VAR、极值尾部拼接、天气状态隐马尔可夫模型，
也没有把 SSP 解释为概率权重。负荷与所有新能源共享一个随机 renewable multiplier，不能
表达 PV、风、其他新能源之间的独立相关矩阵。
\end{gapnote}

\subsection{理论--实现对应}
\begin{center}\small
\begin{tabular}{@{}L{5.0cm}L{8.7cm}@{}}
\toprule
理论条目 & 实现函数与偏差\\
\midrule
二元平稳 AR(1) & \srcpath{scenario\_generation.cpp:correlated\_noise\_factors}\\
块插值与 multiplier clip & 同一函数；插值后逐时相关不等于结点参数\\
确定性 load/PV/wind/other profile & \srcpath{synthesize\_base\_load}、
\srcpath{synthesize\_base\_pv}、\srcpath{synthesize\_base\_wind}、
\srcpath{synthesize\_base\_other\_renewable}\\
储能聚合 SOC & \srcpath{make\_storage\_soc\_profiles}；时域内常值\\
气候查找与形变 & \srcpath{load\_climate\_morph\_factors}、
\srcpath{morph\_regular\_load\_profile}、\srcpath{morph\_regular\_pv}\\
完整多变量天气生成器 & \implnone\\
\bottomrule
\end{tabular}
\end{center}
